\Chapter{92}

\Verse{1}
{
    \zR{话说宝玉从潇湘馆出来，连忙问秋纹道：“老爷叫我作什么？”}
}
{
    \eR{[English source not present in the checked-in Bencraft/Joly files.]}
}
{
}

\Verse{2}
{
    \zR{秋纹笑道：“没 有叫。}
}
{
    \eR{[English source not present in the checked-in Bencraft/Joly files.]}
}
{
}

\Verse{3}
{
    \zR{袭人姐姐叫我请二爷，我怕你不来，才哄你的。”}
}
{
    \eR{[English source not present in the checked-in Bencraft/Joly files.]}
}
{
}

\Verse{4}
{
    \zR{宝玉听了，才把心放下， 因说：“你们请我也罢了，何苦来唬我？”}
}
{
    \eR{[English source not present in the checked-in Bencraft/Joly files.]}
}
{
}

\Verse{5}
{
    \zR{说着，回到怡红院内。}
}
{
    \eR{[English source not present in the checked-in Bencraft/Joly files.]}
}
{
}

\Verse{6}
{
    \zR{袭人便问道：“你 这好半天到那里去了？”}
}
{
    \eR{[English source not present in the checked-in Bencraft/Joly files.]}
}
{
}

\Verse{7}
{
    \zR{宝玉道：“在林姑娘那边，说起姨妈家宝姐姐的事来，就 坐住了。”}
}
{
    \eR{[English source not present in the checked-in Bencraft/Joly files.]}
}
{
}

\Verse{8}
{
    \zR{袭人又问道：“说些什么？”}
}
{
    \eR{[English source not present in the checked-in Bencraft/Joly files.]}
}
{
}

\Verse{9}
{
    \zR{宝玉将打禅语的话述了一遍。}
}
{
    \eR{[English source not present in the checked-in Bencraft/Joly files.]}
}
{
}

\Verse{10}
{
    \zR{袭人道：“你 们再没个计较。}
}
{
    \eR{[English source not present in the checked-in Bencraft/Joly files.]}
}
{
}

\Verse{11}
{
    \zR{正经说些家常闲话儿，或讲究些诗句，也是好的，怎么又说到禅语 上了？}
}
{
    \eR{[English source not present in the checked-in Bencraft/Joly files.]}
}
{
}

\Verse{12}
{
    \zR{又不是和尚。”}
}
{
    \eR{[English source not present in the checked-in Bencraft/Joly files.]}
}
{
}

\Verse{13}
{
    \zR{宝玉道：“你不知道，我们有我们的禅机，别人是插不下嘴 去的。”}
}
{
    \eR{[English source not present in the checked-in Bencraft/Joly files.]}
}
{
}

\Verse{14}
{
    \zR{袭人笑道：“你们参禅参翻了，又叫我们跟着打闷葫芦了。”}
}
{
    \eR{[English source not present in the checked-in Bencraft/Joly files.]}
}
{
}

\Verse{15}
{
    \zR{宝玉道：“头 里我也年纪小，他也孩子气，所以我说了不留神的话，他就恼了。}
}
{
    \eR{[English source not present in the checked-in Bencraft/Joly files.]}
}
{
}

\Verse{16}
{
    \zR{如今我也留神， 他也没有恼的了。}
}
{
    \eR{[English source not present in the checked-in Bencraft/Joly files.]}
}
{
}

\Verse{17}
{
    \zR{只是他近来不常过来，我又念书，偶然到一处，好像生疏了似的。”}
}
{
    \eR{[English source not present in the checked-in Bencraft/Joly files.]}
}
{
}

\Verse{18}
{
    \zR{袭人道：“原该这么着才是。}
}
{
    \eR{[English source not present in the checked-in Bencraft/Joly files.]}
}
{
}

\Verse{19}
{
    \zR{都长了几岁年纪了，怎么好意思还像小孩子时候的样 子？”}
}
{
    \eR{[English source not present in the checked-in Bencraft/Joly files.]}
}
{
}

\Verse{20}
{
    \zR{宝玉点头道：“我也知道。}
}
{
    \eR{[English source not present in the checked-in Bencraft/Joly files.]}
}
{
}

\Verse{21}
{
    \zR{如今且不用说那个。}
}
{
    \eR{[English source not present in the checked-in Bencraft/Joly files.]}
}
{
}

\Verse{22}
{
    \zR{我问你：老太太那里打发人来 说什么来着没有？”}
}
{
    \eR{[English source not present in the checked-in Bencraft/Joly files.]}
}
{
}

\Verse{23}
{
    \zR{袭人道：“没有说什么。”}
}
{
    \eR{[English source not present in the checked-in Bencraft/Joly files.]}
}
{
}

\Verse{24}
{
    \zR{宝玉道：“必是老太太忘了。}
}
{
    \eR{[English source not present in the checked-in Bencraft/Joly files.]}
}
{
}

\Verse{25}
{
    \zR{明儿 不是十一月初一日么？}
}
{
    \eR{[English source not present in the checked-in Bencraft/Joly files.]}
}
{
}

\Verse{26}
{
    \zR{年年老太太那里必是个老规矩，要办消寒会，齐打伙儿坐下 喝酒说笑。}
}
{
    \eR{[English source not present in the checked-in Bencraft/Joly files.]}
}
{
}

\Verse{27}
{
    \zR{我今日已经在学房里告了假了。}
}
{
    \eR{[English source not present in the checked-in Bencraft/Joly files.]}
}
{
}

\Verse{28}
{
    \zR{这会子没有信儿，明儿可是去不去呢？}
}
{
    \eR{[English source not present in the checked-in Bencraft/Joly files.]}
}
{
}

\Verse{29}
{
    \zR{若去了呢，白白的告了假；若不去，老爷知道了，又说我偷懒。”}
}
{
    \eR{[English source not present in the checked-in Bencraft/Joly files.]}
}
{
}

\Verse{30}
{
    \zR{袭人道：“据我 说，你竟是去的是。}
}
{
    \eR{[English source not present in the checked-in Bencraft/Joly files.]}
}
{
}

\Verse{31}
{
    \zR{才念的好些儿了，又想歇着。}
}
{
    \eR{[English source not present in the checked-in Bencraft/Joly files.]}
}
{
}

\Verse{32}
{
    \zR{我劝你也该上点紧儿了。}
}
{
    \eR{[English source not present in the checked-in Bencraft/Joly files.]}
}
{
}

\Verse{33}
{
    \zR{昨儿听 见太太说，兰哥儿念书真好，他打学房里回来，还各自念书作文章，天天晚上弄到 四更多天才睡。}
}
{
    \eR{[English source not present in the checked-in Bencraft/Joly files.]}
}
{
}

\Verse{34}
{
    \zR{你比他大多了，又是叔叔，倘或赶不上他，又叫老太太生气。}
}
{
    \eR{[English source not present in the checked-in Bencraft/Joly files.]}
}
{
}

\Verse{35}
{
    \zR{倒不 如明儿早起去罢。”}
}
{
    \eR{[English source not present in the checked-in Bencraft/Joly files.]}
}
{
}

\Verse{36}
{
    \zR{麝月道：“这么冷天，已经告了假，又去，叫学房里说既这么 着就不该告假呀，显见的是告谎假脱滑儿。}
}
{
    \eR{[English source not present in the checked-in Bencraft/Joly files.]}
}
{
}

\Verse{37}
{
    \zR{依我说，乐得歇一天。}
}
{
    \eR{[English source not present in the checked-in Bencraft/Joly files.]}
}
{
}

\Verse{38}
{
    \zR{就是老太太忘记 了，咱们这里就不消寒了么？}
}
{
    \eR{[English source not present in the checked-in Bencraft/Joly files.]}
}
{
}

\Verse{39}
{
    \zR{咱们也闹个会儿，不好么？”}
}
{
    \eR{[English source not present in the checked-in Bencraft/Joly files.]}
}
{
}

\Verse{40}
{
    \zR{袭人道：“都是你起头 儿，二爷更不肯去了。”}
}
{
    \eR{[English source not present in the checked-in Bencraft/Joly files.]}
}
{
}

\Verse{41}
{
    \zR{麝月道：“我也是乐一天是一天，比不得你要好名儿，使 唤一个月，再多得二两银子。”}
}
{
    \eR{[English source not present in the checked-in Bencraft/Joly files.]}
}
{
}

\Verse{42}
{
    \zR{袭人啐道：“小蹄子儿，人家说正经话，你又来胡 拉混扯的了。”}
}
{
    \eR{[English source not present in the checked-in Bencraft/Joly files.]}
}
{
}

\Verse{43}
{
    \zR{麝月道：“我倒不是混拉扯，我是为你。”}
}
{
    \eR{[English source not present in the checked-in Bencraft/Joly files.]}
}
{
}

\Verse{44}
{
    \zR{袭人道：“为我什么？”}
}
{
    \eR{[English source not present in the checked-in Bencraft/Joly files.]}
}
{
}

\Verse{45}
{
    \zR{麝月道：“二爷上学去了，你又该咕嘟着嘴想着，巴不得二爷早些儿回来，就有说 有笑的了。}
}
{
    \eR{[English source not present in the checked-in Bencraft/Joly files.]}
}
{
}

\Verse{46}
{
    \zR{这会子又假撇清，何苦呢!我都看见了。”}
}
{
    \eR{[English source not present in the checked-in Bencraft/Joly files.]}
}
{
}

\Verse{47}
{
    \zR{袭人正要骂他，只见老太太那里打发人来，说道：“老太太说了，叫二爷明儿 不用上学去呢。}
}
{
    \eR{[English source not present in the checked-in Bencraft/Joly files.]}
}
{
}

\Verse{48}
{
    \zR{明儿请了姨太太来给他解闷，只怕姑娘们都来家里的。}
}
{
    \eR{[English source not present in the checked-in Bencraft/Joly files.]}
}
{
}

\Verse{49}
{
    \zR{史姑娘、邢 姑娘、李姑娘们都请了，明儿来赴什么消寒会呢。”}
}
{
    \eR{[English source not present in the checked-in Bencraft/Joly files.]}
}
{
}

\Verse{50}
{
    \zR{宝玉没有听完，便喜欢道：“可 不是？}
}
{
    \eR{[English source not present in the checked-in Bencraft/Joly files.]}
}
{
}

\Verse{51}
{
    \zR{老太太最高兴的。}
}
{
    \eR{[English source not present in the checked-in Bencraft/Joly files.]}
}
{
}

\Verse{52}
{
    \zR{明日不上学，是过了明路的了。”}
}
{
    \eR{[English source not present in the checked-in Bencraft/Joly files.]}
}
{
}

\Verse{53}
{
    \zR{袭人也不便言语了。}
}
{
    \eR{[English source not present in the checked-in Bencraft/Joly files.]}
}
{
}

\Verse{54}
{
    \zR{那 丫头回去。}
}
{
    \eR{[English source not present in the checked-in Bencraft/Joly files.]}
}
{
}

\Verse{55}
{
    \zR{宝玉认真念了几天书，巴不得玩这一天，又听见薛姨妈过来，想着宝姐 姐自然也来，心里喜欢。}
}
{
    \eR{[English source not present in the checked-in Bencraft/Joly files.]}
}
{
}

\Verse{56}
{
    \zR{便说：“快睡罢，明日早些起来。”}
}
{
    \eR{[English source not present in the checked-in Bencraft/Joly files.]}
}
{
}

\Verse{57}
{
    \zR{于是一夜无话。}
}
{
    \eR{[English source not present in the checked-in Bencraft/Joly files.]}
}
{
}

\Verse{58}
{
    \zR{到了次日，果然一早到老太太那里请了安。}
}
{
    \eR{[English source not present in the checked-in Bencraft/Joly files.]}
}
{
}

\Verse{59}
{
    \zR{又到贾政王夫人那里请了安，回明 了老太太今儿不叫上学，贾政也没言语，便慢慢退出来。}
}
{
    \eR{[English source not present in the checked-in Bencraft/Joly files.]}
}
{
}

\Verse{60}
{
    \zR{走了几步，便一溜烟跑到 贾母房中。}
}
{
    \eR{[English source not present in the checked-in Bencraft/Joly files.]}
}
{
}

\Verse{61}
{
    \zR{见众人都没来，只有凤姐那边的奶妈子，带了巧姐儿，跟着几个小丫头 过来，给老太太请了安，说：“我妈妈先叫我来请安，陪着老太太说说话儿。}
}
{
    \eR{[English source not present in the checked-in Bencraft/Joly files.]}
}
{
}

\Verse{62}
{
    \zR{妈妈 回来就来。”}
}
{
    \eR{[English source not present in the checked-in Bencraft/Joly files.]}
}
{
}

\Verse{63}
{
    \zR{贾母笑着道：“好孩子，我一早就起来了，等他们总不来。}
}
{
    \eR{[English source not present in the checked-in Bencraft/Joly files.]}
}
{
}

\Verse{64}
{
    \zR{只有你二 叔叔来了。”}
}
{
    \eR{[English source not present in the checked-in Bencraft/Joly files.]}
}
{
}

\Verse{65}
{
    \zR{那奶妈子便说：“姑娘，给叔叔请安。”}
}
{
    \eR{[English source not present in the checked-in Bencraft/Joly files.]}
}
{
}

\Verse{66}
{
    \zR{巧姐便请了安。}
}
{
    \eR{[English source not present in the checked-in Bencraft/Joly files.]}
}
{
}

\Verse{67}
{
    \zR{宝玉也问了 一声“妞妞好？”}
}
{
    \eR{[English source not present in the checked-in Bencraft/Joly files.]}
}
{
}

\Verse{68}
{
    \zR{巧姐道：“昨夜听见我妈妈说，要请二叔叔去说话。”}
}
{
    \eR{[English source not present in the checked-in Bencraft/Joly files.]}
}
{
}

\Verse{69}
{
    \zR{宝玉道： “说什么？”}
}
{
    \eR{[English source not present in the checked-in Bencraft/Joly files.]}
}
{
}

\Verse{70}
{
    \zR{巧姐道：“我妈妈说，跟着李妈认了几年字，不知道我认得不认得。}
}
{
    \eR{[English source not present in the checked-in Bencraft/Joly files.]}
}
{
}

\Verse{71}
{
    \zR{我说都认得。}
}
{
    \eR{[English source not present in the checked-in Bencraft/Joly files.]}
}
{
}

\Verse{72}
{
    \zR{我认给妈妈瞧，妈妈说我瞎认，不信，说我一天尽子玩，那里认得。}
}
{
    \eR{[English source not present in the checked-in Bencraft/Joly files.]}
}
{
}

\Verse{73}
{
    \zR{我瞧着那些字也不要紧，就是那《女孝经》也是容易念的。}
}
{
    \eR{[English source not present in the checked-in Bencraft/Joly files.]}
}
{
}

\Verse{74}
{
    \zR{妈妈说我哄他，要请二 叔叔得空儿的时候给我理理。”}
}
{
    \eR{[English source not present in the checked-in Bencraft/Joly files.]}
}
{
}

\Verse{75}
{
    \zR{贾母听了，笑道：“好孩子，你妈妈是不认得字的， 所以说你哄他。}
}
{
    \eR{[English source not present in the checked-in Bencraft/Joly files.]}
}
{
}

\Verse{76}
{
    \zR{明儿叫你二叔叔理给他瞧瞧他就信了。”}
}
{
    \eR{[English source not present in the checked-in Bencraft/Joly files.]}
}
{
}

\Verse{77}
{
    \zR{宝玉道：“你认了多少字 了？”}
}
{
    \eR{[English source not present in the checked-in Bencraft/Joly files.]}
}
{
}

\Verse{78}
{
    \zR{巧姐儿道：“认了二三千多字，念了一本《女孝经》，半个月头里又上了《列 女传》。”}
}
{
    \eR{[English source not present in the checked-in Bencraft/Joly files.]}
}
{
}

\Verse{79}
{
    \zR{宝玉道：“你念了懂的吗？}
}
{
    \eR{[English source not present in the checked-in Bencraft/Joly files.]}
}
{
}

\Verse{80}
{
    \zR{你要不懂，我倒是讲讲这个你听罢。”}
}
{
    \eR{[English source not present in the checked-in Bencraft/Joly files.]}
}
{
}

\Verse{81}
{
    \zR{贾母 道：“做叔叔的也该讲给侄女儿听听。”}
}
{
    \eR{[English source not present in the checked-in Bencraft/Joly files.]}
}
{
}

\Verse{82}
{
    \zR{宝玉便道：“那文王后妃不必说了。}
}
{
    \eR{[English source not present in the checked-in Bencraft/Joly files.]}
}
{
}

\Verse{83}
{
    \zR{那姜后脱簪待罪和齐国的无盐安邦定国， 是后妃里头的贤能的。”}
}
{
    \eR{[English source not present in the checked-in Bencraft/Joly files.]}
}
{
}

\Verse{84}
{
    \zR{巧姐听了，答应个“是”。}
}
{
    \eR{[English source not present in the checked-in Bencraft/Joly files.]}
}
{
}

\Verse{85}
{
    \zR{宝玉又道：“若说有才的，是 曹大姑、班婕妤、蔡文姬、谢道韫诸人。”}
}
{
    \eR{[English source not present in the checked-in Bencraft/Joly files.]}
}
{
}

\Verse{86}
{
    \zR{巧姐问道：“那贤德的呢？”}
}
{
    \eR{[English source not present in the checked-in Bencraft/Joly files.]}
}
{
}

\Verse{87}
{
    \zR{宝玉道： “孟光的荆钗布裙，鲍宣妻的提瓮出汲，陶侃母的截发留宾：这些不厌贫的，就是 贤德了。”}
}
{
    \eR{[English source not present in the checked-in Bencraft/Joly files.]}
}
{
}

\Verse{88}
{
    \zR{巧姐欣然点头。}
}
{
    \eR{[English source not present in the checked-in Bencraft/Joly files.]}
}
{
}

\Verse{89}
{
    \zR{宝玉道：“还有苦的，像那乐昌破镜，苏蕙回文；那孝 的，木兰代父从军，曹娥投水寻尸等类，也难尽说。”}
}
{
    \eR{[English source not present in the checked-in Bencraft/Joly files.]}
}
{
}

\Verse{90}
{
    \zR{巧姐听到这些，却默默如有 所思。}
}
{
    \eR{[English source not present in the checked-in Bencraft/Joly files.]}
}
{
}

\Verse{91}
{
    \zR{宝玉又讲那曹氏的引刀割鼻及那些守节的，巧姐听着更觉肃敬起来。}
}
{
    \eR{[English source not present in the checked-in Bencraft/Joly files.]}
}
{
}

\Verse{92}
{
    \zR{宝玉恐 他不自在，又说：“那些艳的，如王嫱、西子、樊素、小蛮、绛仙、文君、红拂， 都是女中的――”尚未说出，贾母见巧姐默然，便说：“够了，不用说了。}
}
{
    \eR{[English source not present in the checked-in Bencraft/Joly files.]}
}
{
}

\Verse{93}
{
    \zR{讲的太 多，他那里记得。”}
}
{
    \eR{[English source not present in the checked-in Bencraft/Joly files.]}
}
{
}

\Verse{94}
{
    \zR{巧姐道：“二叔叔才说的，也有念过的，也有没念过的。}
}
{
    \eR{[English source not present in the checked-in Bencraft/Joly files.]}
}
{
}

\Verse{95}
{
    \zR{念过 的一讲我更知道好处了。”}
}
{
    \eR{[English source not present in the checked-in Bencraft/Joly files.]}
}
{
}

\Verse{96}
{
    \zR{宝玉道：“那字是自然认得的，不用再理了。”}
}
{
    \eR{[English source not present in the checked-in Bencraft/Joly files.]}
}
{
}

\Verse{97}
{
    \zR{巧姐道：“我还听见我妈妈说：我们家的小红，头里是二叔叔那里的，我妈妈 要了来，还没有补上人呢。}
}
{
    \eR{[English source not present in the checked-in Bencraft/Joly files.]}
}
{
}

\Verse{98}
{
    \zR{我妈妈想着要把什么柳家的五儿补上，不知二叔叔要不 要。”}
}
{
    \eR{[English source not present in the checked-in Bencraft/Joly files.]}
}
{
}

\Verse{99}
{
    \zR{宝玉听了更喜欢，笑着道：“你听你妈妈的话!要补谁就补谁罢咧，又问什 么要不要呢。”}
}
{
    \eR{[English source not present in the checked-in Bencraft/Joly files.]}
}
{
}

\Verse{100}
{
    \zR{因又向贾母笑道：“我瞧大妞妞这个小模样儿，又有这个聪明儿， 只怕将来比凤姐姐还强呢，又比他认的字。”}
}
{
    \eR{[English source not present in the checked-in Bencraft/Joly files.]}
}
{
}

\Verse{101}
{
    \zR{贾母道：“女孩儿家认得字也好，只 是女工针黹倒是要紧的。”}
}
{
    \eR{[English source not present in the checked-in Bencraft/Joly files.]}
}
{
}

\Verse{102}
{
    \zR{巧姐儿道：“我也跟着刘妈妈学着做呢。}
}
{
    \eR{[English source not present in the checked-in Bencraft/Joly files.]}
}
{
}

\Verse{103}
{
    \zR{什么扎花儿咧， 拉锁子咧，我虽弄不好，却也学着会做几针儿。”}
}
{
    \eR{[English source not present in the checked-in Bencraft/Joly files.]}
}
{
}

\Verse{104}
{
    \zR{贾母道：“咱们这样人家，固然 不仗着自己做，但只到底知道些，日后才不受人家的拿捏。”}
}
{
    \eR{[English source not present in the checked-in Bencraft/Joly files.]}
}
{
}

\Verse{105}
{
    \zR{巧姐答应着“是”， 还要宝玉解说《列女传》，见宝玉呆呆的，也不好再问。}
}
{
    \eR{[English source not present in the checked-in Bencraft/Joly files.]}
}
{
}

\Verse{106}
{
    \zR{你道宝玉呆的是什么？}
}
{
    \eR{[English source not present in the checked-in Bencraft/Joly files.]}
}
{
}

\Verse{107}
{
    \zR{只 因柳五儿要进怡红院，头一次是他病了，不能进来，第二次王夫人撵了晴雯，大凡 有些姿色的，都不敢挑。}
}
{
    \eR{[English source not present in the checked-in Bencraft/Joly files.]}
}
{
}

\Verse{108}
{
    \zR{后来又在吴贵家看晴雯去，五儿跟着他妈给晴雯送东西去， 见了一面，更觉娇娜妩媚。}
}
{
    \eR{[English source not present in the checked-in Bencraft/Joly files.]}
}
{
}

\Verse{109}
{
    \zR{今日亏得凤姐想着，叫他补入小红的窝儿，竟是喜出望 外了，所以呆呆的呆想。}
}
{
    \eR{[English source not present in the checked-in Bencraft/Joly files.]}
}
{
}

\Verse{110}
{
    \zR{贾母等着那些人，见这时候还不来，又叫丫头去请。}
}
{
    \eR{[English source not present in the checked-in Bencraft/Joly files.]}
}
{
}

\Verse{111}
{
    \zR{回来李纨同着他妹子、探 春、惜春、史湘云、黛玉都来了。}
}
{
    \eR{[English source not present in the checked-in Bencraft/Joly files.]}
}
{
}

\Verse{112}
{
    \zR{大家请了贾母的安，众人厮见。}
}
{
    \eR{[English source not present in the checked-in Bencraft/Joly files.]}
}
{
}

\Verse{113}
{
    \zR{独有薛姨妈未到， 贾母又叫请去。}
}
{
    \eR{[English source not present in the checked-in Bencraft/Joly files.]}
}
{
}

\Verse{114}
{
    \zR{果然薛姨妈带着宝琴过来。}
}
{
    \eR{[English source not present in the checked-in Bencraft/Joly files.]}
}
{
}

\Verse{115}
{
    \zR{宝玉请了安，问了好，只不见宝钗邢岫 烟二人。}
}
{
    \eR{[English source not present in the checked-in Bencraft/Joly files.]}
}
{
}

\Verse{116}
{
    \zR{黛玉便问起：“宝姐姐为何不来？”}
}
{
    \eR{[English source not present in the checked-in Bencraft/Joly files.]}
}
{
}

\Verse{117}
{
    \zR{薛姨妈假说身上不好。}
}
{
    \eR{[English source not present in the checked-in Bencraft/Joly files.]}
}
{
}

\Verse{118}
{
    \zR{邢岫烟知道薛 姨妈在坐，所以不来。}
}
{
    \eR{[English source not present in the checked-in Bencraft/Joly files.]}
}
{
}

\Verse{119}
{
    \zR{宝玉虽见宝钗不来，心中纳闷，因黛玉来了，便把想宝钗的 心暂且搁开。}
}
{
    \eR{[English source not present in the checked-in Bencraft/Joly files.]}
}
{
}

\Verse{120}
{
    \zR{不多时，邢王二夫人也来了。}
}
{
    \eR{[English source not present in the checked-in Bencraft/Joly files.]}
}
{
}

\Verse{121}
{
    \zR{凤姐听见婆婆们先到了，自己不好落后， 只得打发平儿先来告假，说是：“正要过来，因身上发热，过一回儿就来。”}
}
{
    \eR{[English source not present in the checked-in Bencraft/Joly files.]}
}
{
}

\Verse{122}
{
    \zR{贾母 道：“既是身上不好，不来也罢。}
}
{
    \eR{[English source not present in the checked-in Bencraft/Joly files.]}
}
{
}

\Verse{123}
{
    \zR{咱们这时候很该吃饭了。”}
}
{
    \eR{[English source not present in the checked-in Bencraft/Joly files.]}
}
{
}

\Verse{124}
{
    \zR{丫头们把火盆往后挪 了一挪，就在贾母榻前一溜摆下两桌，大家序次坐下。}
}
{
    \eR{[English source not present in the checked-in Bencraft/Joly files.]}
}
{
}

\Verse{125}
{
    \zR{吃了饭，依旧围炉闲谈，不 须多赘。}
}
{
    \eR{[English source not present in the checked-in Bencraft/Joly files.]}
}
{
}

\Verse{126}
{
    \zR{且说凤姐因何不来？}
}
{
    \eR{[English source not present in the checked-in Bencraft/Joly files.]}
}
{
}

\Verse{127}
{
    \zR{头里为着倒比邢王二夫人迟了不好意思，后来旺儿家的来 回说：“迎姑娘那里打发人来请奶奶安，还说并没有到上头，只到奶奶这里来。”}
}
{
    \eR{[English source not present in the checked-in Bencraft/Joly files.]}
}
{
}

\Verse{128}
{
    \zR{凤姐听了纳闷，不知又是什么事，便叫那人进来，问：“姑娘在家好？”}
}
{
    \eR{[English source not present in the checked-in Bencraft/Joly files.]}
}
{
}

\Verse{129}
{
    \zR{那人道： “有什么好的。}
}
{
    \eR{[English source not present in the checked-in Bencraft/Joly files.]}
}
{
}

\Verse{130}
{
    \zR{奴才并不是姑娘打发来的，实在是司棋的母亲要我来求奶奶的。”}
}
{
    \eR{[English source not present in the checked-in Bencraft/Joly files.]}
}
{
}

\Verse{131}
{
    \zR{凤姐道：“司棋已经出去了，为什么来求我？”}
}
{
    \eR{[English source not present in the checked-in Bencraft/Joly files.]}
}
{
}

\Verse{132}
{
    \zR{那人道：“自从司棋出去，终日啼 哭。}
}
{
    \eR{[English source not present in the checked-in Bencraft/Joly files.]}
}
{
}

\Verse{133}
{
    \zR{忽然那一日，他表兄来了。}
}
{
    \eR{[English source not present in the checked-in Bencraft/Joly files.]}
}
{
}

\Verse{134}
{
    \zR{他母亲见了，恨的什么儿似的，说他害了司棋，一 把拉住要打。}
}
{
    \eR{[English source not present in the checked-in Bencraft/Joly files.]}
}
{
}

\Verse{135}
{
    \zR{那小子不敢言语。}
}
{
    \eR{[English source not present in the checked-in Bencraft/Joly files.]}
}
{
}

\Verse{136}
{
    \zR{谁知司棋听见了，急忙出来，老着脸，和他母亲说： ‘我是为他出来的，我也恨他没良心。}
}
{
    \eR{[English source not present in the checked-in Bencraft/Joly files.]}
}
{
}

\Verse{137}
{
    \zR{如今他来了，妈要打他，不如勒死了我罢。’}
}
{
    \eR{[English source not present in the checked-in Bencraft/Joly files.]}
}
{
}

\Verse{138}
{
    \zR{他妈骂他：‘不害臊的东西，你心里要怎么样？’}
}
{
    \eR{[English source not present in the checked-in Bencraft/Joly files.]}
}
{
}

\Verse{139}
{
    \zR{司棋说道：‘一个女人嫁一个男 人。}
}
{
    \eR{[English source not present in the checked-in Bencraft/Joly files.]}
}
{
}

\Verse{140}
{
    \zR{我一时失脚，上了他的当，我就是他的人了，决不肯再跟着别人的。}
}
{
    \eR{[English source not present in the checked-in Bencraft/Joly files.]}
}
{
}

\Verse{141}
{
    \zR{我只恨他 为什么这么胆小，一身作事一身当，为什么逃了呢？}
}
{
    \eR{[English source not present in the checked-in Bencraft/Joly files.]}
}
{
}

\Verse{142}
{
    \zR{就是他一辈子不来，我也一辈 子不嫁人的。}
}
{
    \eR{[English source not present in the checked-in Bencraft/Joly files.]}
}
{
}

\Verse{143}
{
    \zR{妈要给我配人，我原拚着一死。}
}
{
    \eR{[English source not present in the checked-in Bencraft/Joly files.]}
}
{
}

\Verse{144}
{
    \zR{今儿他来了，妈问他怎么样。}
}
{
    \eR{[English source not present in the checked-in Bencraft/Joly files.]}
}
{
}

\Verse{145}
{
    \zR{要是他 不改心，我在妈跟前磕了头，只当是我死了，他到那里，我跟到那里，就是讨饭吃 也是愿意的。’}
}
{
    \eR{[English source not present in the checked-in Bencraft/Joly files.]}
}
{
}

\Verse{146}
{
    \zR{他妈气的了不得，便哭着骂着说：‘你是我的女儿，我偏不给他， 你敢怎么着？’}
}
{
    \eR{[English source not present in the checked-in Bencraft/Joly files.]}
}
{
}

\Verse{147}
{
    \zR{那知道司棋这东西糊涂，便一头撞在墙上，把脑袋撞破，鲜血流出， 竟碰死了。}
}
{
    \eR{[English source not present in the checked-in Bencraft/Joly files.]}
}
{
}

\Verse{148}
{
    \zR{他妈哭着，救不过来，便要叫那小子偿命。}
}
{
    \eR{[English source not present in the checked-in Bencraft/Joly files.]}
}
{
}

\Verse{149}
{
    \zR{他表兄也奇，说道：‘你们 不用着急。}
}
{
    \eR{[English source not present in the checked-in Bencraft/Joly files.]}
}
{
}

\Verse{150}
{
    \zR{我在外头原发了财，因想着他才回来的，心也算是真了。}
}
{
    \eR{[English source not present in the checked-in Bencraft/Joly files.]}
}
{
}

\Verse{151}
{
    \zR{你们要不信， 只管瞧。’}
}
{
    \eR{[English source not present in the checked-in Bencraft/Joly files.]}
}
{
}

\Verse{152}
{
    \zR{说着，打怀里掏出一匣子金珠首饰来。}
}
{
    \eR{[English source not present in the checked-in Bencraft/Joly files.]}
}
{
}

\Verse{153}
{
    \zR{他妈妈看见了，心软了，说：‘你 既有心，为什么总不言语？’}
}
{
    \eR{[English source not present in the checked-in Bencraft/Joly files.]}
}
{
}

\Verse{154}
{
    \zR{他外甥道：‘大凡女人都是水性杨花，我要说有钱， 他就是贪图银钱了。}
}
{
    \eR{[English source not present in the checked-in Bencraft/Joly files.]}
}
{
}

\Verse{155}
{
    \zR{如今他这为人就是难得的。}
}
{
    \eR{[English source not present in the checked-in Bencraft/Joly files.]}
}
{
}

\Verse{156}
{
    \zR{我把首饰给你们，我去买棺盛殓 他。’}
}
{
    \eR{[English source not present in the checked-in Bencraft/Joly files.]}
}
{
}

\Verse{157}
{
    \zR{那司棋的母亲接了东西，也不顾女孩儿了，由着外甥去。}
}
{
    \eR{[English source not present in the checked-in Bencraft/Joly files.]}
}
{
}

\Verse{158}
{
    \zR{那里知道他外甥叫 人抬了两口棺材来。}
}
{
    \eR{[English source not present in the checked-in Bencraft/Joly files.]}
}
{
}

\Verse{159}
{
    \zR{司棋的母亲看见诧异，说怎么棺材要两口，他外甥笑道：‘一 口装不下，得两口才好。’}
}
{
    \eR{[English source not present in the checked-in Bencraft/Joly files.]}
}
{
}

\Verse{160}
{
    \zR{司棋的母亲见他外甥又不哭，只当是他心疼的傻了。}
}
{
    \eR{[English source not present in the checked-in Bencraft/Joly files.]}
}
{
}

\Verse{161}
{
    \zR{岂 知他忙着把司棋收拾了，也不啼哭，眼错不见，把带的小刀子往脖子里一抹，也就 抹死了。}
}
{
    \eR{[English source not present in the checked-in Bencraft/Joly files.]}
}
{
}

\Verse{162}
{
    \zR{司棋的母亲懊悔起来，倒哭的了不得。}
}
{
    \eR{[English source not present in the checked-in Bencraft/Joly files.]}
}
{
}

\Verse{163}
{
    \zR{如今坊里知道了，要报官。}
}
{
    \eR{[English source not present in the checked-in Bencraft/Joly files.]}
}
{
}

\Verse{164}
{
    \zR{他急了， 央我来求奶奶说个人情，他再过来给奶奶磕头。”}
}
{
    \eR{[English source not present in the checked-in Bencraft/Joly files.]}
}
{
}

\Verse{165}
{
    \zR{凤姐听了，诧异道：“那有这样傻丫头，偏偏的就碰见这个傻小子!怪不得那 一天翻出那些东西来，他心里没事人似的，敢只是这么个烈性孩子。}
}
{
    \eR{[English source not present in the checked-in Bencraft/Joly files.]}
}
{
}

\Verse{166}
{
    \zR{论起来我也没 这么大工夫管他这些闲事，但只你才说的，叫人听着怪可怜见儿的。}
}
{
    \eR{[English source not present in the checked-in Bencraft/Joly files.]}
}
{
}

\Verse{167}
{
    \zR{也罢了，你回 去告诉他，我和你二爷说，打发旺儿给他撕掳就是了。”}
}
{
    \eR{[English source not present in the checked-in Bencraft/Joly files.]}
}
{
}

\Verse{168}
{
    \zR{凤姐打发那人去了，才过 贾母这边来，不提。}
}
{
    \eR{[English source not present in the checked-in Bencraft/Joly files.]}
}
{
}

\Verse{169}
{
    \zR{且说贾政这日正与詹光下大棋，通局的输赢也差不多，单为着一只角儿死活未 分，在那里打结。}
}
{
    \eR{[English source not present in the checked-in Bencraft/Joly files.]}
}
{
}

\Verse{170}
{
    \zR{门上的小厮进来回道：“外面冯大爷要见老爷。”}
}
{
    \eR{[English source not present in the checked-in Bencraft/Joly files.]}
}
{
}

\Verse{171}
{
    \zR{贾政说：“请 进来。”}
}
{
    \eR{[English source not present in the checked-in Bencraft/Joly files.]}
}
{
}

\Verse{172}
{
    \zR{小厮出去请了，冯紫英走进门来，贾政即忙迎着。}
}
{
    \eR{[English source not present in the checked-in Bencraft/Joly files.]}
}
{
}

\Verse{173}
{
    \zR{冯紫英进来，在书房中 坐下，见是下棋，便道：“只管下棋，我来观局。”}
}
{
    \eR{[English source not present in the checked-in Bencraft/Joly files.]}
}
{
}

\Verse{174}
{
    \zR{詹光笑道：“晚生的棋是不堪 瞧的。”}
}
{
    \eR{[English source not present in the checked-in Bencraft/Joly files.]}
}
{
}

\Verse{175}
{
    \zR{冯紫英道：“好说，请下罢。”}
}
{
    \eR{[English source not present in the checked-in Bencraft/Joly files.]}
}
{
}

\Verse{176}
{
    \zR{贾政道：“有什么事么？”}
}
{
    \eR{[English source not present in the checked-in Bencraft/Joly files.]}
}
{
}

\Verse{177}
{
    \zR{冯紫英道：“没 有什么话。}
}
{
    \eR{[English source not present in the checked-in Bencraft/Joly files.]}
}
{
}

\Verse{178}
{
    \zR{老伯只管下棋，我也学几着儿。”}
}
{
    \eR{[English source not present in the checked-in Bencraft/Joly files.]}
}
{
}

\Verse{179}
{
    \zR{贾政向詹光道：“冯大爷是我们相好 的，既没事，我们索性下完了这一局再说话儿。}
}
{
    \eR{[English source not present in the checked-in Bencraft/Joly files.]}
}
{
}

\Verse{180}
{
    \zR{冯大爷在旁边瞧着。”}
}
{
    \eR{[English source not present in the checked-in Bencraft/Joly files.]}
}
{
}

\Verse{181}
{
    \zR{冯紫英道： “下采不下采？”}
}
{
    \eR{[English source not present in the checked-in Bencraft/Joly files.]}
}
{
}

\Verse{182}
{
    \zR{詹光道：“下采的。”}
}
{
    \eR{[English source not present in the checked-in Bencraft/Joly files.]}
}
{
}

\Verse{183}
{
    \zR{冯紫英道：“下采的是不好多嘴的。”}
}
{
    \eR{[English source not present in the checked-in Bencraft/Joly files.]}
}
{
}

\Verse{184}
{
    \zR{贾 政道：“多嘴也不妨，横竖他输了十来两银子，终久是不拿出来的。}
}
{
    \eR{[English source not present in the checked-in Bencraft/Joly files.]}
}
{
}

\Verse{185}
{
    \zR{往后只好罚他 做东便了。”}
}
{
    \eR{[English source not present in the checked-in Bencraft/Joly files.]}
}
{
}

\Verse{186}
{
    \zR{詹光笑道：“这倒使得。”}
}
{
    \eR{[English source not present in the checked-in Bencraft/Joly files.]}
}
{
}

\Verse{187}
{
    \zR{冯紫英道：“老伯和詹公对下么？”}
}
{
    \eR{[English source not present in the checked-in Bencraft/Joly files.]}
}
{
}

\Verse{188}
{
    \zR{贾政 笑道：“从前对下，他输了；如今让他两个子儿，他又输了。}
}
{
    \eR{[English source not present in the checked-in Bencraft/Joly files.]}
}
{
}

\Verse{189}
{
    \zR{时常还要悔几着，不 叫他悔他就急了。”}
}
{
    \eR{[English source not present in the checked-in Bencraft/Joly files.]}
}
{
}

\Verse{190}
{
    \zR{詹光也笑道：“没有的事。”}
}
{
    \eR{[English source not present in the checked-in Bencraft/Joly files.]}
}
{
}

\Verse{191}
{
    \zR{贾政道：“你试试瞧。”}
}
{
    \eR{[English source not present in the checked-in Bencraft/Joly files.]}
}
{
}

\Verse{192}
{
    \zR{大家一 面说笑，一面下完了。}
}
{
    \eR{[English source not present in the checked-in Bencraft/Joly files.]}
}
{
}

\Verse{193}
{
    \zR{做起棋来，詹光还了棋头，输了七个子儿。}
}
{
    \eR{[English source not present in the checked-in Bencraft/Joly files.]}
}
{
}

\Verse{194}
{
    \zR{冯紫英道：“这 盘总吃亏在打结里头。}
}
{
    \eR{[English source not present in the checked-in Bencraft/Joly files.]}
}
{
}

\Verse{195}
{
    \zR{老伯结少，就便宜了。”}
}
{
    \eR{[English source not present in the checked-in Bencraft/Joly files.]}
}
{
}

\Verse{196}
{
    \zR{贾政对冯紫英道：“有罪，有罪，咱们说话儿罢。”}
}
{
    \eR{[English source not present in the checked-in Bencraft/Joly files.]}
}
{
}

\Verse{197}
{
    \zR{冯紫英道：“小侄与老伯 久不见面。}
}
{
    \eR{[English source not present in the checked-in Bencraft/Joly files.]}
}
{
}

\Verse{198}
{
    \zR{一来会会，二来因广西的同知进来引见，带了四种洋货，可以做得贡的。}
}
{
    \eR{[English source not present in the checked-in Bencraft/Joly files.]}
}
{
}

\Verse{199}
{
    \zR{一件是围屏，有二十四扇？}
}
{
    \eR{[English source not present in the checked-in Bencraft/Joly files.]}
}
{
}

\Verse{200}
{
    \zR{子，都是紫檀雕刻的。}
}
{
    \eR{[English source not present in the checked-in Bencraft/Joly files.]}
}
{
}

\Verse{201}
{
    \zR{中间虽说不是玉，却是绝好的硝 子石，石上镂出山水、人物、楼台、花鸟儿来。}
}
{
    \eR{[English source not present in the checked-in Bencraft/Joly files.]}
}
{
}

\Verse{202}
{
    \zR{一扇上有五六十个人，都是宫妆的 女子，名为‘汉宫春晓’。}
}
{
    \eR{[English source not present in the checked-in Bencraft/Joly files.]}
}
{
}

\Verse{203}
{
    \zR{人的眉、目、口、鼻以及出手、衣褶，刻得又清楚，又 细腻。}
}
{
    \eR{[English source not present in the checked-in Bencraft/Joly files.]}
}
{
}

\Verse{204}
{
    \zR{点缀布置，都是好的。}
}
{
    \eR{[English source not present in the checked-in Bencraft/Joly files.]}
}
{
}

\Verse{205}
{
    \zR{我想尊府大观园中正厅上恰好用的着。}
}
{
    \eR{[English source not present in the checked-in Bencraft/Joly files.]}
}
{
}

\Verse{206}
{
    \zR{还有一架钟表， 有三尺多高，也是一个童儿拿着时辰牌，到什么时候儿就报什么时辰。}
}
{
    \eR{[English source not present in the checked-in Bencraft/Joly files.]}
}
{
}

\Verse{207}
{
    \zR{里头还有消 息人儿打十番儿。}
}
{
    \eR{[English source not present in the checked-in Bencraft/Joly files.]}
}
{
}

\Verse{208}
{
    \zR{这是两件重笨的，却还没有拿来。}
}
{
    \eR{[English source not present in the checked-in Bencraft/Joly files.]}
}
{
}

\Verse{209}
{
    \zR{现在我带在这里的两件，却倒 有些意思儿。”}
}
{
    \eR{[English source not present in the checked-in Bencraft/Joly files.]}
}
{
}

\Verse{210}
{
    \zR{就在身边拿出一个锦匣子来，用几重白绫裹着。}
}
{
    \eR{[English source not present in the checked-in Bencraft/Joly files.]}
}
{
}

\Verse{211}
{
    \zR{揭开了绵子，第一 层是一个玻璃盒子，里头金托子大红绉绸托底，上放着一颗桂圆大的珠子，光华耀 目。}
}
{
    \eR{[English source not present in the checked-in Bencraft/Joly files.]}
}
{
}

\Verse{212}
{
    \zR{冯紫英道：“据说这就叫做‘母珠’。”}
}
{
    \eR{[English source not present in the checked-in Bencraft/Joly files.]}
}
{
}

\Verse{213}
{
    \zR{因叫：“拿一个盘儿来。”}
}
{
    \eR{[English source not present in the checked-in Bencraft/Joly files.]}
}
{
}

\Verse{214}
{
    \zR{詹光即忙 端过一个黑漆茶盘，道：“使得么？”}
}
{
    \eR{[English source not present in the checked-in Bencraft/Joly files.]}
}
{
}

\Verse{215}
{
    \zR{冯紫英道：“使得。”}
}
{
    \eR{[English source not present in the checked-in Bencraft/Joly files.]}
}
{
}

\Verse{216}
{
    \zR{便又向怀里掏出一个 白绢包儿，将包儿里的珠子都倒在盘里散着，把那颗母珠搁在中间，将盘放于桌上。}
}
{
    \eR{[English source not present in the checked-in Bencraft/Joly files.]}
}
{
}

\Verse{217}
{
    \zR{看见那些小珠子儿滴溜滴溜的都滚到大珠子身边，回来把这颗大珠子抬高了，别处 的小珠子一颗也不剩，都粘在大珠上。}
}
{
    \eR{[English source not present in the checked-in Bencraft/Joly files.]}
}
{
}

\Verse{218}
{
    \zR{詹光道：“这也奇！”}
}
{
    \eR{[English source not present in the checked-in Bencraft/Joly files.]}
}
{
}

\Verse{219}
{
    \zR{贾政道：“这是有的， 所以叫做‘母珠’，原是珠之母。”}
}
{
    \eR{[English source not present in the checked-in Bencraft/Joly files.]}
}
{
}

\Verse{220}
{
    \zR{那冯紫英又回头看着他跟来的小厮道：“那个匣子呢？”}
}
{
    \eR{[English source not present in the checked-in Bencraft/Joly files.]}
}
{
}

\Verse{221}
{
    \zR{小厮赶忙捧过一个花 梨木匣子来。}
}
{
    \eR{[English source not present in the checked-in Bencraft/Joly files.]}
}
{
}

\Verse{222}
{
    \zR{大家打开看时，原来匣内衬着虎纹锦，锦上叠着一束蓝纱。}
}
{
    \eR{[English source not present in the checked-in Bencraft/Joly files.]}
}
{
}

\Verse{223}
{
    \zR{詹光道： “这是什么东西？”}
}
{
    \eR{[English source not present in the checked-in Bencraft/Joly files.]}
}
{
}

\Verse{224}
{
    \zR{冯紫英道：“这叫做‘鲛绡帐’。”}
}
{
    \eR{[English source not present in the checked-in Bencraft/Joly files.]}
}
{
}

\Verse{225}
{
    \zR{在匣子里拿出来时，叠得 长不满五寸，厚不上半寸。}
}
{
    \eR{[English source not present in the checked-in Bencraft/Joly files.]}
}
{
}

\Verse{226}
{
    \zR{冯紫英一层一层的打开，打到十来层，已经桌上铺不下 了。}
}
{
    \eR{[English source not present in the checked-in Bencraft/Joly files.]}
}
{
}

\Verse{227}
{
    \zR{冯紫英道：“你看，里头还有两褶，必得高屋里去才张得下。}
}
{
    \eR{[English source not present in the checked-in Bencraft/Joly files.]}
}
{
}

\Verse{228}
{
    \zR{这就是鲛丝所织。}
}
{
    \eR{[English source not present in the checked-in Bencraft/Joly files.]}
}
{
}

\Verse{229}
{
    \zR{暑热天气张在堂屋里头，苍蝇蚊子一个不能进来，又轻又亮。”}
}
{
    \eR{[English source not present in the checked-in Bencraft/Joly files.]}
}
{
}

\Verse{230}
{
    \zR{贾政道：“不用全 打开，怕叠起来倒费事。”}
}
{
    \eR{[English source not present in the checked-in Bencraft/Joly files.]}
}
{
}

\Verse{231}
{
    \zR{詹光便与冯紫英一层一层折好收拾了。}
}
{
    \eR{[English source not present in the checked-in Bencraft/Joly files.]}
}
{
}

\Verse{232}
{
    \zR{冯紫英道：“这四件东西，价儿也不贵，两万银他就卖。}
}
{
    \eR{[English source not present in the checked-in Bencraft/Joly files.]}
}
{
}

\Verse{233}
{
    \zR{母珠一万，鲛绡帐五 千，‘汉宫春晓’与自鸣钟五千。”}
}
{
    \eR{[English source not present in the checked-in Bencraft/Joly files.]}
}
{
}

\Verse{234}
{
    \zR{贾政道：“那里买的起！”}
}
{
    \eR{[English source not present in the checked-in Bencraft/Joly files.]}
}
{
}

\Verse{235}
{
    \zR{冯紫英道：“你们 是个国戚，难道宫里头用不着么？”}
}
{
    \eR{[English source not present in the checked-in Bencraft/Joly files.]}
}
{
}

\Verse{236}
{
    \zR{贾政道：“用得着的很多，只是那里有这些银 子？}
}
{
    \eR{[English source not present in the checked-in Bencraft/Joly files.]}
}
{
}

\Verse{237}
{
    \zR{等我叫人拿进去给老太太瞧瞧。”}
}
{
    \eR{[English source not present in the checked-in Bencraft/Joly files.]}
}
{
}

\Verse{238}
{
    \zR{冯紫英道：“很是。”}
}
{
    \eR{[English source not present in the checked-in Bencraft/Joly files.]}
}
{
}

\Verse{239}
{
    \zR{贾政便着人叫贾琏把这两件东西送到老太太那边去，并叫人请了邢王二夫人、 凤姐儿都来瞧着，又把两件东西一一试过。}
}
{
    \eR{[English source not present in the checked-in Bencraft/Joly files.]}
}
{
}

\Verse{240}
{
    \zR{贾琏道：“他还有两件：一件是围屏， 一件是乐钟。}
}
{
    \eR{[English source not present in the checked-in Bencraft/Joly files.]}
}
{
}

\Verse{241}
{
    \zR{共总要卖二万银子呢。”}
}
{
    \eR{[English source not present in the checked-in Bencraft/Joly files.]}
}
{
}

\Verse{242}
{
    \zR{凤姐儿接着道：“东西自然是好的，但是那 里有这些闲钱？}
}
{
    \eR{[English source not present in the checked-in Bencraft/Joly files.]}
}
{
}

\Verse{243}
{
    \zR{咱们又不比外任督抚要办贡。}
}
{
    \eR{[English source not present in the checked-in Bencraft/Joly files.]}
}
{
}

\Verse{244}
{
    \zR{我已经想了好些年了，像咱们这种人 家，必得置些不动摇的根基才好：或是祭地，或是义庄，再置些坟屋。}
}
{
    \eR{[English source not present in the checked-in Bencraft/Joly files.]}
}
{
}

\Verse{245}
{
    \zR{往后子孙遇 见不得意的事，还是点儿底子，不到一败涂地。}
}
{
    \eR{[English source not present in the checked-in Bencraft/Joly files.]}
}
{
}

\Verse{246}
{
    \zR{我的意思是这样，不知老太太、老 爷、太太们怎么样？}
}
{
    \eR{[English source not present in the checked-in Bencraft/Joly files.]}
}
{
}

\Verse{247}
{
    \zR{若是外头老爷们要买只管买。”}
}
{
    \eR{[English source not present in the checked-in Bencraft/Joly files.]}
}
{
}

\Verse{248}
{
    \zR{贾母与众人都说：“这话说的 倒也是。”}
}
{
    \eR{[English source not present in the checked-in Bencraft/Joly files.]}
}
{
}

\Verse{249}
{
    \zR{贾琏道：“还了他罢。}
}
{
    \eR{[English source not present in the checked-in Bencraft/Joly files.]}
}
{
}

\Verse{250}
{
    \zR{原是老爷叫我送给老太太瞧，为的是宫里好进， 谁说买来搁在家里？}
}
{
    \eR{[English source not present in the checked-in Bencraft/Joly files.]}
}
{
}

\Verse{251}
{
    \zR{老太太还没开口，你便说了一大堆丧气话。”}
}
{
    \eR{[English source not present in the checked-in Bencraft/Joly files.]}
}
{
}

\Verse{252}
{
    \zR{说着，便把两件 东西拿出去了，告诉贾政，只说：“老太太不要。”}
}
{
    \eR{[English source not present in the checked-in Bencraft/Joly files.]}
}
{
}

\Verse{253}
{
    \zR{便与冯紫英道：“这两件东西 好可好，就只没银子。}
}
{
    \eR{[English source not present in the checked-in Bencraft/Joly files.]}
}
{
}

\Verse{254}
{
    \zR{我替你留心，有要买的人我便送信给你去。”}
}
{
    \eR{[English source not present in the checked-in Bencraft/Joly files.]}
}
{
}

\Verse{255}
{
    \zR{冯紫英只得收 拾好了，坐下说些闲话，没有兴头，就要起身。}
}
{
    \eR{[English source not present in the checked-in Bencraft/Joly files.]}
}
{
}

\Verse{256}
{
    \zR{贾政道：“你在这里吃了晚饭去罢。”}
}
{
    \eR{[English source not present in the checked-in Bencraft/Joly files.]}
}
{
}

\Verse{257}
{
    \zR{冯紫英道：“罢了，来了就叨搅老伯吗？”}
}
{
    \eR{[English source not present in the checked-in Bencraft/Joly files.]}
}
{
}

\Verse{258}
{
    \zR{贾政道：“说那里的话。”}
}
{
    \eR{[English source not present in the checked-in Bencraft/Joly files.]}
}
{
}

\Verse{259}
{
    \zR{正说着，人回：“大老爷来了。”}
}
{
    \eR{[English source not present in the checked-in Bencraft/Joly files.]}
}
{
}

\Verse{260}
{
    \zR{贾赦早已进来。}
}
{
    \eR{[English source not present in the checked-in Bencraft/Joly files.]}
}
{
}

\Verse{261}
{
    \zR{彼此相见，叙些寒温。}
}
{
    \eR{[English source not present in the checked-in Bencraft/Joly files.]}
}
{
}

\Verse{262}
{
    \zR{不一 时摆上酒来，肴馔罗列，大家喝着酒。}
}
{
    \eR{[English source not present in the checked-in Bencraft/Joly files.]}
}
{
}

\Verse{263}
{
    \zR{至四五巡后，说起洋货的话。}
}
{
    \eR{[English source not present in the checked-in Bencraft/Joly files.]}
}
{
}

\Verse{264}
{
    \zR{冯紫英道：“这 种货本是难消的。}
}
{
    \eR{[English source not present in the checked-in Bencraft/Joly files.]}
}
{
}

\Verse{265}
{
    \zR{除非要像尊府这样人家还可消得，其馀就难了。”}
}
{
    \eR{[English source not present in the checked-in Bencraft/Joly files.]}
}
{
}

\Verse{266}
{
    \zR{贾政道：“这 也不见得。”}
}
{
    \eR{[English source not present in the checked-in Bencraft/Joly files.]}
}
{
}

\Verse{267}
{
    \zR{贾赦道：“我们家里也比不得从前了，这回儿也不过是个空门面。”}
}
{
    \eR{[English source not present in the checked-in Bencraft/Joly files.]}
}
{
}

\Verse{268}
{
    \zR{冯紫英又问：“东府珍大爷可好么？}
}
{
    \eR{[English source not present in the checked-in Bencraft/Joly files.]}
}
{
}

\Verse{269}
{
    \zR{我前儿见他，说起家常话儿来，提到他令郎续 娶的媳妇远不及头里那位秦氏奶奶了。}
}
{
    \eR{[English source not present in the checked-in Bencraft/Joly files.]}
}
{
}

\Verse{270}
{
    \zR{如今后娶的到底是那一家的？}
}
{
    \eR{[English source not present in the checked-in Bencraft/Joly files.]}
}
{
}

\Verse{271}
{
    \zR{我也没有问 起。”}
}
{
    \eR{[English source not present in the checked-in Bencraft/Joly files.]}
}
{
}

\Verse{272}
{
    \zR{贾政道：“我们这个侄孙媳妇儿也是这里大家，从前做过京畿道的胡老爷的 女孩儿。”}
}
{
    \eR{[English source not present in the checked-in Bencraft/Joly files.]}
}
{
}

\Verse{273}
{
    \zR{冯紫英道：“胡道长我是知道的，但是他家教上也不怎么样。}
}
{
    \eR{[English source not present in the checked-in Bencraft/Joly files.]}
}
{
}

\Verse{274}
{
    \zR{也罢了， 只要姑娘好就好。”}
}
{
    \eR{[English source not present in the checked-in Bencraft/Joly files.]}
}
{
}

\Verse{275}
{
    \zR{贾琏道：“听得内阁里人说起，雨村又要升了。”}
}
{
    \eR{[English source not present in the checked-in Bencraft/Joly files.]}
}
{
}

\Verse{276}
{
    \zR{贾政道：“这也好。}
}
{
    \eR{[English source not present in the checked-in Bencraft/Joly files.]}
}
{
}

\Verse{277}
{
    \zR{不知准 不准？”}
}
{
    \eR{[English source not present in the checked-in Bencraft/Joly files.]}
}
{
}

\Verse{278}
{
    \zR{贾琏道：“大约有意思的了。”}
}
{
    \eR{[English source not present in the checked-in Bencraft/Joly files.]}
}
{
}

\Verse{279}
{
    \zR{冯紫英道：“我今儿从吏部里来，也听见 这样说。}
}
{
    \eR{[English source not present in the checked-in Bencraft/Joly files.]}
}
{
}

\Verse{280}
{
    \zR{雨村老先生是贵本家不是？”}
}
{
    \eR{[English source not present in the checked-in Bencraft/Joly files.]}
}
{
}

\Verse{281}
{
    \zR{贾政道：“是。”}
}
{
    \eR{[English source not present in the checked-in Bencraft/Joly files.]}
}
{
}

\Verse{282}
{
    \zR{冯紫英道：“是有服的， 还是无服的？”}
}
{
    \eR{[English source not present in the checked-in Bencraft/Joly files.]}
}
{
}

\Verse{283}
{
    \zR{贾政道：“说也话长。}
}
{
    \eR{[English source not present in the checked-in Bencraft/Joly files.]}
}
{
}

\Verse{284}
{
    \zR{他原籍是浙江湖州府人，流寓到苏州，甚不 得意。}
}
{
    \eR{[English source not present in the checked-in Bencraft/Joly files.]}
}
{
}

\Verse{285}
{
    \zR{有个甄士隐和他相好，时常周济他。}
}
{
    \eR{[English source not present in the checked-in Bencraft/Joly files.]}
}
{
}

\Verse{286}
{
    \zR{以后中了进士，得了榜下知县，便娶了 甄家的丫头。}
}
{
    \eR{[English source not present in the checked-in Bencraft/Joly files.]}
}
{
}

\Verse{287}
{
    \zR{如今的太太不是正配。}
}
{
    \eR{[English source not present in the checked-in Bencraft/Joly files.]}
}
{
}

\Verse{288}
{
    \zR{岂知甄士隐弄到零落不堪，没有找处。}
}
{
    \eR{[English source not present in the checked-in Bencraft/Joly files.]}
}
{
}

\Verse{289}
{
    \zR{雨村革 了职以后，那时还与我家并未相识，只因舍妹丈林如海林公在扬州巡盐的时候，请 他在家做西席，外甥女儿是他的学生。}
}
{
    \eR{[English source not present in the checked-in Bencraft/Joly files.]}
}
{
}

\Verse{290}
{
    \zR{因他有起复的信，要进京来，恰好外甥女儿 要上来探亲，林姑老爷便托他照应上来的，还有一封荐书托我吹嘘吹嘘。}
}
{
    \eR{[English source not present in the checked-in Bencraft/Joly files.]}
}
{
}

\Verse{291}
{
    \zR{那时看他 不错，大家常会。}
}
{
    \eR{[English source not present in the checked-in Bencraft/Joly files.]}
}
{
}

\Verse{292}
{
    \zR{岂知雨村也奇：我家世袭起，从‘代’字辈下来，宁荣两宅，人 口房舍，以及起居事宜，一概都明白。}
}
{
    \eR{[English source not present in the checked-in Bencraft/Joly files.]}
}
{
}

\Verse{293}
{
    \zR{因此，遂觉得亲热了。”}
}
{
    \eR{[English source not present in the checked-in Bencraft/Joly files.]}
}
{
}

\Verse{294}
{
    \zR{因又笑说道：“几 年间，门子也会钻了，由知府推升转了御史，不过几年，升了吏部侍郎，兵部尚书。}
}
{
    \eR{[English source not present in the checked-in Bencraft/Joly files.]}
}
{
}

\Verse{295}
{
    \zR{为着一件事降了三级，如今又要升了。”}
}
{
    \eR{[English source not present in the checked-in Bencraft/Joly files.]}
}
{
}

\Verse{296}
{
    \zR{冯紫英道：“人世的荣枯，仕途的得失，终属难定。”}
}
{
    \eR{[English source not present in the checked-in Bencraft/Joly files.]}
}
{
}

\Verse{297}
{
    \zR{贾政道：“天下事都是 一个样的理哟。}
}
{
    \eR{[English source not present in the checked-in Bencraft/Joly files.]}
}
{
}

\Verse{298}
{
    \zR{比如方才那珠子，那颗大的就像有福气的人似的，那些小的都托赖 着他的灵气护庇着。}
}
{
    \eR{[English source not present in the checked-in Bencraft/Joly files.]}
}
{
}

\Verse{299}
{
    \zR{要是那大的没有了，那些小的也就没有收揽了。}
}
{
    \eR{[English source not present in the checked-in Bencraft/Joly files.]}
}
{
}

\Verse{300}
{
    \zR{就像人家儿当 头人有了事，骨肉也都分离了，亲戚也都零落了，就是好朋友也都散了。}
}
{
    \eR{[English source not present in the checked-in Bencraft/Joly files.]}
}
{
}

\Verse{301}
{
    \zR{转瞬荣枯， 真似春云秋叶一般。}
}
{
    \eR{[English source not present in the checked-in Bencraft/Joly files.]}
}
{
}

\Verse{302}
{
    \zR{你想做官有什么趣儿呢？}
}
{
    \eR{[English source not present in the checked-in Bencraft/Joly files.]}
}
{
}

\Verse{303}
{
    \zR{像雨村算便宜的了。}
}
{
    \eR{[English source not present in the checked-in Bencraft/Joly files.]}
}
{
}

\Verse{304}
{
    \zR{还有我们差不多 的人家儿，就是甄家，从前一样功勋，一样世袭，一样起居，我们也是时常来往。}
}
{
    \eR{[English source not present in the checked-in Bencraft/Joly files.]}
}
{
}

\Verse{305}
{
    \zR{不多几年他们进京来，差人到我这里请安，还很热闹。}
}
{
    \eR{[English source not present in the checked-in Bencraft/Joly files.]}
}
{
}

\Verse{306}
{
    \zR{一会儿抄了原籍的家财，至 今杳无音信。}
}
{
    \eR{[English source not present in the checked-in Bencraft/Joly files.]}
}
{
}

\Verse{307}
{
    \zR{不知他近况若何，心下也着实惦记着。”}
}
{
    \eR{[English source not present in the checked-in Bencraft/Joly files.]}
}
{
}

\Verse{308}
{
    \zR{贾赦道：“什么珠子？”}
}
{
    \eR{[English source not present in the checked-in Bencraft/Joly files.]}
}
{
}

\Verse{309}
{
    \zR{贾 政同冯紫英又说了一遍给贾赦听。}
}
{
    \eR{[English source not present in the checked-in Bencraft/Joly files.]}
}
{
}

\Verse{310}
{
    \zR{贾赦道：“咱们家是再没有事的。”}
}
{
    \eR{[English source not present in the checked-in Bencraft/Joly files.]}
}
{
}

\Verse{311}
{
    \zR{冯紫英道： “果然尊府是不怕的。}
}
{
    \eR{[English source not present in the checked-in Bencraft/Joly files.]}
}
{
}

\Verse{312}
{
    \zR{一则里头有贵妃照应；二则故旧好，亲戚多；三则你们家自 老太太起，至于少爷们，没有一个刁钻刻薄的。”}
}
{
    \eR{[English source not present in the checked-in Bencraft/Joly files.]}
}
{
}

\Verse{313}
{
    \zR{贾政道：“虽无刁钻刻薄的，却 没有德行才情。}
}
{
    \eR{[English source not present in the checked-in Bencraft/Joly files.]}
}
{
}

\Verse{314}
{
    \zR{白白的衣租食税，那里当得起？”}
}
{
    \eR{[English source not present in the checked-in Bencraft/Joly files.]}
}
{
}

\Verse{315}
{
    \zR{贾赦道：“咱们不用说这些话， 大家吃酒罢。”}
}
{
    \eR{[English source not present in the checked-in Bencraft/Joly files.]}
}
{
}

\Verse{316}
{
    \zR{大家又喝了几杯，摆上饭来。}
}
{
    \eR{[English source not present in the checked-in Bencraft/Joly files.]}
}
{
}

\Verse{317}
{
    \zR{吃毕喝茶，冯家的小厮走来，轻轻的向紫英说了 一句。}
}
{
    \eR{[English source not present in the checked-in Bencraft/Joly files.]}
}
{
}

\Verse{318}
{
    \zR{冯紫英便要告辞。}
}
{
    \eR{[English source not present in the checked-in Bencraft/Joly files.]}
}
{
}

\Verse{319}
{
    \zR{贾赦问那小厮道：“你说什么？”}
}
{
    \eR{[English source not present in the checked-in Bencraft/Joly files.]}
}
{
}

\Verse{320}
{
    \zR{小厮道：“外面下雪， 早已下了梆子了。”}
}
{
    \eR{[English source not present in the checked-in Bencraft/Joly files.]}
}
{
}

\Verse{321}
{
    \zR{贾政叫人看时，已是雪深一寸多了。}
}
{
    \eR{[English source not present in the checked-in Bencraft/Joly files.]}
}
{
}

\Verse{322}
{
    \zR{贾政道：“那两件东西， 你收拾好了么？”}
}
{
    \eR{[English source not present in the checked-in Bencraft/Joly files.]}
}
{
}

\Verse{323}
{
    \zR{冯紫英道：“收好了。}
}
{
    \eR{[English source not present in the checked-in Bencraft/Joly files.]}
}
{
}

\Verse{324}
{
    \zR{若尊府要用，价钱还自然让些。”}
}
{
    \eR{[English source not present in the checked-in Bencraft/Joly files.]}
}
{
}

\Verse{325}
{
    \zR{贾政道： “我留神就是了。”}
}
{
    \eR{[English source not present in the checked-in Bencraft/Joly files.]}
}
{
}

\Verse{326}
{
    \zR{紫英道：“我再听信罢。}
}
{
    \eR{[English source not present in the checked-in Bencraft/Joly files.]}
}
{
}

\Verse{327}
{
    \zR{天气冷，请罢，别送了。”}
}
{
    \eR{[English source not present in the checked-in Bencraft/Joly files.]}
}
{
}

\Verse{328}
{
    \zR{贾赦贾政 便命贾琏送了出去。}
}
{
    \eR{[English source not present in the checked-in Bencraft/Joly files.]}
}
{
}

\Verse{329}
{
    \zR{未知后事如何，下回分解。}
}
{
    \eR{[English source not present in the checked-in Bencraft/Joly files.]}
}
{
}

\Verse{330}
{
    \zR{(本章完)}
}
{
    \eR{[English source not present in the checked-in Bencraft/Joly files.]}
}
{
}
